\DocumentMetadata{
  pdfversion=2.0,
  pdfstandard=ua-2,
  lang=en-US,
  tagging=on,
  tagging-setup={math/setup=mathml-SE,tabsorder=structure}
}
\documentclass[11pt,letterpaper]{article}
\usepackage[margin=0.68in,includefoot]{geometry}
\usepackage{newcomputermodern}
\usepackage{amsmath,amscd,mathtools}
\usepackage{tikz,adjustbox}
\providecommand{\square}{\mdlgwhtsquare}
\usepackage[hidelinks]{hyperref}
\hypersetup{
  pdftitle={Analysis First-Year Exam, May 2023, Part 2},
  pdfauthor={University of Florida Department of Mathematics},
  pdfsubject={Graduate mathematics examination},
  pdfdisplaydoctitle=true,
  bookmarksopen=true
}
\setcounter{secnumdepth}{0}
\setlength{\parindent}{0pt}
\setlength{\parskip}{0.28em}
\setlength{\emergencystretch}{3em}
\setlength{\leftmargini}{2.15em}
\setlength{\leftmarginii}{2.65em}
\setlength{\leftmarginiii}{3em}
\pagestyle{empty}
\raggedbottom
\allowdisplaybreaks
\NewTaggingSocketPlug{block/list/label}{examproblem}{%
  \tagstructbegin{tag=Lbl}\tagstructend%
  \tagstructbegin{tag=\LBody}%
  \tagstructbegin{tag=\ProblemHeadingTag,title-o={\ProblemHeadingText}}%
  \tagmcbegin{tag=\ProblemHeadingTag,actualtext={\ProblemHeadingText}}%
  #2%
  \tagmcend%
  \tagstructend%
}
\newenvironment{examproblems}[2]{%
  \def\ProblemHeadingTag{#2}%
  \AssignTaggingSocketPlug{block/list/label}{examproblem}%
  \begin{enumerate}%
  \setcounter{enumi}{#1}%
  \renewcommand{\labelenumi}{\textbf{\theenumi.}}%
  \setlength{\topsep}{0.35em}%
  \setlength{\partopsep}{0pt}%
  \setlength{\itemsep}{0.48em}%
  \setlength{\parsep}{0.18em}%
}{\end{enumerate}}
\newenvironment{examsubparts}{%
  \AssignTaggingSocketPlug{block/list/label}{default}%
  \begin{enumerate}%
  \setlength{\topsep}{0.18em}%
  \setlength{\partopsep}{0pt}%
  \setlength{\itemsep}{0.16em}%
  \setlength{\parsep}{0.1em}%
}{\end{enumerate}}
\newenvironment{examinstructions}{%
  \par\begingroup\itshape
}{%
  \par\endgroup\vspace{0.18em}
}
\makeatletter
\renewcommand\subsection{\@startsection{subsection}{2}{0pt}%
  {-1.25ex plus -.25ex minus -.1ex}%
  {0.55ex plus .1ex}%
  {\normalfont\large\bfseries}}
\renewcommand{\@maketitle}{%
  \newpage\null\vskip -1em%
  {\footnotesize\bfseries
    \href{https://www.ufl.edu/}{University of Florida}, \href{https://math.ufl.edu/}{Department of Mathematics}\par}%
  \vskip -0.5em%
  \tagstructbegin{tag=H1,title={Analysis First-Year Exam, May 2023, Part 2}}%
  \tagpdfparaOff%
  \tagmcbegin{tag=H1}%
    {\LARGE\bfseries\href{https://gma.math.ufl.edu/exams/analysis-first-year/}{Analysis First-Year Exam}, May 2023, Part 2\par}%
  \tagmcend%
  \tagpdfparaOn%
  \tagstructend%
  \vskip 0.25em%
  \tagmcbegin{artifact}%
    \hrule height 1pt%
  \tagmcend%
  \vskip 1em%
}
\makeatother
\title{Analysis First-Year Exam, May 2023, Part 2}
\author{}
\date{}
\begin{document}
\maketitle
\thispagestyle{empty}
\begin{examinstructions}
Answer FOUR questions in detail. State carefully any results used without proof.
\end{examinstructions}
\begin{examproblems}{0}{H2}
\def\ProblemHeadingText{Problem 1.}
\item
Prove that the function \(f:[a,b]\to\mathbb{R}\) is Riemann-integrable if and only if for each \(\epsilon>0\) there exist Riemann-integrable functions~\(\ell\) and~\(u\) on \([a,b]\) such that \(\ell\leq f\leq u\) and \(\int_a^b(u-\ell)<\epsilon\).
\def\ProblemHeadingText{Problem 2.}
\item
Let \(f:[0,1]\to\mathbb{R}\) be continuous; for each \(n\in\mathbb{N}\) and each \(t\in[0,1]\) define \(f_n(t)=f(t)t^n\). Prove that the sequence \((f_n:n\in\mathbb{N})\) converges uniformly on \([0,1]\) if and only if \(f(1)\) has a specific value, which should be found.
\def\ProblemHeadingText{Problem 3.}
\item
Let \(f:[-1,1]\to\mathbb{R}\) be continuous. Prove that if
\[
\int_{-1}^1 f(t)t^n\,dt=0
\]
 whenever the natural number~\(n\) is odd, then the function~\(f\) is even.
\def\ProblemHeadingText{Problem 4.}
\item
Let \((f_n:n\in\mathbb{N})\) be a sequence of measurable real-valued functions on some measurable space. Prove that each of these three sets is measurable:
\[
\begin{aligned}
C&=\{\omega:\text{ the sequence }f_n(\omega)\text{ is Cauchy}\};\\
D&=\{\omega:\text{ the sequence }f_n(\omega)\text{ has all terms different}\};\\
E&=\{\omega:\text{ the sequence }f_n(\omega)\text{ is eventually constant}\}.
\end{aligned}
\]
\def\ProblemHeadingText{Problem 5.}
\item
The bounded function \(f:[0,1]\times[0,1]\to\mathbb{R}\) is separately continuous: that is, \(f(x,y)\) is continuous in~\(x\) when~\(y\) is fixed and continuous in~\(y\) when~\(x\) is fixed. Prove the continuity of the function \(F:[0,1]\to\mathbb{R}\) defined by
\[
F(x)=\int_0^1 f(x,y)\,dy.
\]
\end{examproblems}
\end{document}
